\setcounter{section}{22}

Wir besprechen nun eine wichtige analytische Hilfstechnik namens \stichwort {Partition der Eins} {.} Wir werden sie im Beweis für die Aussage, dass orientierbare Mannigfaltigkeiten eine positive Volumenform besitzen, und für den Beweis des Satzes von Stokes einsetzen. In dieser Vorlesung werden wir Partitionen der Eins konstruieren, wozu wir zunächst einige topologische Begriffe benötigen.

  \zwischenueberschrift{Kompakte Ausschöpfung}

\bild{ \begin{center}
\includegraphics[width=5.5cm]{ \bildeinlesungpng {Inner_point} {png} }
\end{center}
\bildtext {Das offene Innere ist die Vereinigung aller inneren Punkte, also derjenigen Punkte von $T$
\zusatzklammer {im Bild $E$} {} {,}
die mit einer ganzen offenen Umgebung in $T$ enthalten sind.} }

\bildlizenz { Inner point.png } {} {Zasdfgbnm} {Commons} {PD} {}

\inputdefinition
{}
{

Es sei $X$ ein
\definitionsverweis {topologischer Raum}{}{}
und

\mavergleichskette
{\vergleichskette
{T
}
{ \subseteq }{ X
}
{  }{
}
{    }{
}
{   }{
}
}
{}{}{}
eine Teilmenge. Dann heißt

\mavergleichskettedisp
{\vergleichskette
{ \overline{  T  }
}
{ =} {  \bigcap_{A \text{ abgeschlossen}, \, T \subseteq A} A
}
{ } {
}
{ } {
}
{ } {
}
}
{}{}{}
der \definitionswort {Abschluss}{}
\zusatzklammer {oder \definitionswort {topologische Abschluss}{}} {} {}
von $T$.

}
Für metrische Räume haben wir den Abschluss als Menge aller
\definitionsverweis {Berührpunkte}{}{}
schon früher eingeführt, siehe insbesondere
[[Metrischer Raum/Abschluss/Berührungspunkte/Durchschnitt/Aufgabe|Aufgabe 35.3 (Analysis (Osnabrück 2021-2023))]].

\inputdefinition
{}
{

Es sei $X$ ein
\definitionsverweis {topologischer Raum}{}{}
und

\mavergleichskette
{\vergleichskette
{ T
}
{ \subseteq }{ X
}
{  }{
}
{    }{
}
{   }{
}
}
{}{}{}
eine Teilmenge. Dann heißt

\mavergleichskettedisp
{\vergleichskette
{ T ^{o}
}
{ =} { \bigcup_{U \text{ offen}, \, U \subseteq T}   U
}
{ } {
}
{ } {
}
{ } {
}
}
{}{}{}
das \definitionswort {offene Innere}{}
\zusatzklammer {oder \stichwort {Innere} {}} {} {}
von $T$.

}
Diese beiden Begriffe stehen durch

\mavergleichskettedisp
{\vergleichskette
{ \overline{ T }
}
{ =} { X \setminus (X \setminus T)^{o}
}
{ } {
}
{ } {
}
{ } {
}
}
{}{}{}
miteinander in Beziehung. Auch der Begriff des Randes überträgt sich von der metrischen Situation auf beliebige topologische Räume.

\inputdefinition
{
}
{

Es sei $X$ ein
\definitionsverweis {topologischer Raum}{}{}
und

\mavergleichskette
{\vergleichskette
{ T
}
{ \subseteq }{ X
}
{  }{
}
{    }{
}
{   }{
}
}
{}{}{}
eine Teilmenge. Unter dem
\definitionswort {Rand}{}
von $T$ versteht man die Menge

\mavergleichskettedisp
{\vergleichskette
{ \operatorname{Rand} \, ( T )
}
{ =} { \overline{  T  }  \setminus T ^{o}
}
{ } {
}
{ } {
}
{ } {
}
}
{}{}{.}

}
Man beachte, dass dieser topologische Rand ein anderes Konzept ist als der Rand bei einer berandeten Mannigfaltigkeit, allerdings besteht eine Beziehung, die in
[[Mannigfaltigkeit mit Rand/Rand ist topologischer Rand des Komplementes/Aufgabe|Aufgabe 22.1]]
besprochen wird.

\inputdefinition
{}
{

Es sei $X$ ein
\definitionsverweis {topologischer Raum}{}{}
und
\maabbdisp {f} { X } {\R
} {}
eine
\definitionsverweis {Funktion}{}{.}
Dann heißt der
\definitionsverweis {topologische Abschluss}{}{}

\mathdisp {\overline{  \left\{  x \in X  \mid  f(x) \neq 0        \right\}   }} {  }

der \definitionswort {Träger}{} von $f$.

}

\inputdefinition
{}
{

Es sei $X$ ein
\definitionsverweis {topologischer Raum}{}{.}
Eine \definitionswort {kompakte Ausschöp\-fung}{}

\mathbed {A_n} {}
 {n \in \N} {}
 {} {} {} {,}
von $X$ ist eine
\definitionsverweis {Folge}{}{}
von
\definitionsverweis {kompakten Teilmengen}{}{}

\mavergleichskette
{\vergleichskette
{ A_n
}
{ \subseteq }{ X
}
{  }{
}
{    }{
}
{   }{
}
}
{}{}{}
mit

\mathdisp {A_n \subseteq A_{n+1}^{o} \text{ und } \bigcup_{n=0}^\infty A_n = X} { . }

}

Der $\R^d$ wird durch die abgeschlossenen Bälle

\mathbed {B  \left( 0,n \right)} {}
 {n \in \N} {}
 {} {} {} {.}
kompakt ausgeschöpft, siehe
[[R^n/Kompakte Ausschöpfung/Aufgabe|Aufgabe 22.6]].

  \zwischenueberschrift{Partitionen der Eins}

__NOEDITSECTION__

\inputfaktbeweis
{Mannigfaltigkeit/Abzählbare Basis/Kompakte Ausschöpfung/Fakt}
{Lemma}
{}
{

\faktsituation {Es sei $M$ eine
\definitionsverweis {Mannigfaltigkeit}{}{}}
\faktvoraussetzung {mit einer
\definitionsverweis {abzählbaren Basis der Topologie}{}{.}}
\faktfolgerung {Dann besitzt $M$ eine
\definitionsverweis {kompakte Ausschöpfung}{}{.}}
\faktzusatz {}
\faktzusatz {}

}
{

Zu jedem Punkt

\mavergleichskette
{\vergleichskette
{ P
}
{ \in }{ M
}
{  }{
}
{    }{
}
{   }{
}
}
{}{}{}
gibt es eine offene Kartenumgebung

\mavergleichskette
{\vergleichskette
{ P
}
{ \in }{ U
}
{  }{
}
{    }{
}
{   }{
}
}
{}{}{,}
\maabbdisp {\alpha} { U } { V
} {}
sowie Ballumgebungen

\mavergleichskettedisp
{\vergleichskette
{ U \left(  \alpha(P), \epsilon  \right)
}
{ \subseteq} { B  \left( \alpha(P), \epsilon \right)
}
{ \subset} { V
}
{ } {
}
{ } {
}
}
{}{}{.}
Wegen der
\definitionsverweis {Homöomorphie}{}{}
der Kartenabbildung und der Kompaktheit der abgeschlossenen Bälle ist

\mavergleichskette
{\vergleichskette
{ B_P
}
{ = }{ \alpha^{-1} \left(  B  \left( \alpha(P), \epsilon \right)  \right)
}
{  }{
}
{    }{
}
{   }{
}
}
{}{}{}
eine
\definitionsverweis {kompakte Teilmenge}{}{}
von $M$, die die
\definitionsverweis {offene Umgebung}{}{}

\mavergleichskette
{\vergleichskette
{ U_P
}
{ = }{ \alpha^{-1} \left(  U \left(  \alpha(P), \epsilon  \right)  \right)
}
{  }{
}
{    }{
}
{   }{
}
}
{}{}{}
von $P$ umfasst. Die

\mathbed {U_P} {}
 {P \in M} {}
 {} {} {} {,}
bilden eine
\definitionsverweis {offene Überdeckung}{}{}
von $M$, sodass es nach
[[Topologischer Raum/Abzählbare Basis/Überdeckung besitzt abzählbare Teilüberdeckung/Aufgabe|Aufgabe 2.8 (Maß- und Integrationstheorie (Osnabrück 2022-2023))]]
eine abzählbare Teilüberdeckung gibt. Diese sei mit

\mathbed {U_n} {}
 {n \in \N} {}
 {} {} {} {,}
bezeichnet
\zusatzklammer {wobei die $U_n$ in den kompakten Teilmengen $B_n$ liegen} {} {.}
Wir definieren nun
\definitionsverweis {rekursiv}{}{}
eine
\definitionsverweis {monoton wachsende}{}{}
Abbildung
\maabbeledisp {} {\N} {\N
} { k } {n_k
} {,}
derart, dass

\mathbeddisp {A_k = \bigcup_{n=0}^{n_k } B_n} {}
 {k \in \N} {}
 {} {} {} {,}
eine
\definitionsverweis {kompakte Ausschöpfung}{}{}
von $M$ ist. Als endliche Vereinigungen von kompakten Mengen sind diese $A_k$ kompakt. Wir beginnen mit

\mavergleichskette
{\vergleichskette
{}
{ n_0 }{}
{  }{}
{    }{}
{   }{}
}
{}{}{.}
Es sei $n_k$ schon konstruiert. Die Menge

\mathdisp {A_k \cup B_{n_k +1}} {  }

ist
\definitionsverweis {kompakt}{}{}
und wird daher von endlich vielen offenen Mengen
\mathl{\bigcup_{n=0}^{n_{k+1} } U_n}{} überdeckt, wobei wir

\mavergleichskette
{\vergleichskette
{ n_{k+1}
}
{ \geq }{ n_k +1
}
{  }{
}
{    }{
}
{   }{
}
}
{}{}{}
wählen. Mit dieser Wahl ist

\mavergleichskettedisp
{\vergleichskette
{ A_k
}
{ \subseteq} { \bigcup_{n = 0}^{n_{k+1} } U_n
}
{ \subseteq} { \bigcup_{n = 0}^{n_{k+1} } B_n
}
{ =} { A_{k+1}
}
{ } {
}
}
{}{}{,}
und diese Folge bildet eine Ausschöpfung, da die

\mathbed {U_n} {}
 {n \in \N} {}
 {} {} {} {,}
eine offene Überdeckung von $M$  bilden.

}

\bild{ \begin{center}
\includegraphics[width=5.5cm]{ \bildeinlesungsvg {Partition_of_unity_illustration} {svg} }
\end{center}
\bildtext {} }

\bildlizenz { Partition of unity illustration.svg } {} {Oleg Alexandrov} {Commons} {PD} {}

\inputdefinition
{}
{

Es sei $X$ ein
\definitionsverweis {topologischer Raum}{}{.}
Eine Familie von Funktionen
\maabbdisp {h_j} { X } {\R
} {}
mit

\mavergleichskette
{\vergleichskette
{ j
}
{ \in }{ J
}
{  }{
}
{    }{
}
{   }{
}
}
{}{}{}
heißt eine \definitionswort {Partition der Eins}{,} wenn folgende Eigenschaften gelten.
\aufzaehlungdrei{Es ist

\mavergleichskette
{\vergleichskette
{ h_j(X)
}
{ \subseteq }{ [0,1]
}
{  }{
}
{    }{
}
{   }{
}
}
{}{}{}
für alle

\mavergleichskette
{\vergleichskette
{ j
}
{ \in }{ J
}
{  }{
}
{    }{
}
{   }{
}
}
{}{}{.}
}{Jeder Punkt

\mavergleichskette
{\vergleichskette
{ P
}
{ \in }{ X
}
{  }{
}
{    }{
}
{   }{
}
}
{}{}{}
besitzt eine
\definitionsverweis {offene Umgebung}{}{}

\mavergleichskette
{\vergleichskette
{ P
}
{ \in }{ U
}
{  }{
}
{    }{
}
{   }{
}
}
{}{}{}
derart, dass die
\definitionsverweis {eingeschränkten Funktionen}{}{}

\mathl{h_j {{|}}_{ U }}{} bis auf endlich viele Ausnahmen die Nullfunktion sind.
}{Es ist

\mavergleichskette
{\vergleichskette
{  \sum_{j \in J} h_j
}
{ = }{ 1
}
{  }{
}
{    }{
}
{   }{
}
}
{}{}{.}
} Wenn alle $h_j$ \definitionsverweis {stetig}{}{} sind, so spricht man von einer \definitionswort {stetigen Partition der Eins}{.}

}

Die zweite Eigenschaft sichert dabei, dass die Summe in (3) definiert ist, da für jeden Punkt

\mavergleichskette
{\vergleichskette
{ P
}
{ \in }{ X
}
{  }{
}
{    }{
}
{   }{
}
}
{}{}{}
und fast alle

\mavergleichskette
{\vergleichskette
{ j
}
{ \in }{ J
}
{  }{
}
{    }{
}
{   }{
}
}
{}{}{}
die Gleichheit

\mavergleichskette
{\vergleichskette
{  h_j(P)
}
{ = }{ 0
}
{  }{
}
{    }{
}
{   }{
}
}
{}{}{}
gilt. Bei einer Mannigfaltigkeit nennt man eine solche Partition \stichwort {differenzierbar} {,} wenn alle $h_j$
\definitionsverweis {differenzierbare Funktionen}{}{}
sind.

\inputdefinition
{}
{

Es sei

\mavergleichskette
{\vergleichskette
{ X
}
{ = }{ \bigcup_{i \in I} W_i
}
{  }{
}
{    }{
}
{   }{
}
}
{}{}{}
eine
\definitionsverweis {offene Überdeckung}{}{}
eines
\definitionsverweis {topologischen Raumes}{}{}
$X$. Eine
\definitionsverweis {Partition der Eins}{}{}
\maabbdisp {h_j} { X } {\R
} {}
mit

\mavergleichskette
{\vergleichskette
{ j
}
{ \in }{ J
}
{  }{
}
{    }{
}
{   }{
}
}
{}{}{}
heißt eine \definitionswort {der Überdeckung untergeordnete Partition der Eins}{,} wenn es für jedes

\mavergleichskette
{\vergleichskette
{ j
}
{ \in }{ J
}
{  }{
}
{    }{
}
{   }{
}
}
{}{}{}
eine offene Menge
\mathl{W_{i(j)}}{} aus der Überdeckung derart gibt, dass der
\definitionsverweis {Träger}{}{}
von $h_j$ in
\mathl{W_{i(j)}}{} liegt.

}

__NOEDITSECTION__

\inputfaktbeweis
{Mannigfaltigkeit/Abzählbare Basis/Lokal endlicher Atlas/Fakt}
{Lemma}
{}
{

\faktsituation {Es sei $M$ eine
\definitionsverweis {Mannigfaltigkeit}{}{}
mit einer
\definitionsverweis {abzählbaren Basis}{}{}
der Topologie. Es sei

\mavergleichskette
{\vergleichskette
{ M
}
{ = }{ \bigcup_{i \in I} W _i
}
{  }{
}
{    }{
}
{   }{
}
}
{}{}{}
eine
\definitionsverweis {offene Überdeckung}{}{}
von $M$.}
\faktfolgerung {Dann gibt es einen abzählbaren
\definitionsverweis {verträglichen Atlas}{}{}

\mathbed {\left(  U_j, \alpha_j, V_j  \right)} {}
 {j \in J} {}
 {} {} {} {,}
mit Ballumgebungen

\mavergleichskettedisp
{\vergleichskette
{ U \left(   0 , \delta_j   \right)
}
{ \subset} { B  \left(  0 , \epsilon_j  \right)
}
{ \subset} { V_j
}
{ } {
}
{ } {
}
}
{}{}{}
\zusatzklammer {dabei ist

\mavergleichskettek
{\vergleichskettek
{ 0
}
{ \in }{ V_j
}
{ \subseteq }{ \R^n
}
{    }{
}
{   }{
}
}
{}{}{}
und

\mavergleichskettek
{\vergleichskettek
{ \delta_j
}
{ < }{ \epsilon_j
}
{  }{
}
{    }{
}
{   }{
}
}
{}{}{}} {} {}
derart, dass es für jedes

\mavergleichskette
{\vergleichskette
{ j
}
{ \in }{ J
}
{  }{
}
{    }{
}
{   }{
}
}
{}{}{}
ein
\mathl{W_{i(j)}}{} mit

\mavergleichskette
{\vergleichskette
{ U_j
}
{ \subseteq }{ W_{i(j)}
}
{  }{
}
{    }{
}
{   }{
}
}
{}{}{}
gibt, dass $M$ von

\mathbed {\alpha_j^{-1} \left(  U \left(   0 , \delta_j   \right)  \right)} {}
 {j \in J} {}
 {} {} {} {,}
überdeckt wird und dass jeder Punkt

\mavergleichskette
{\vergleichskette
{ P
}
{ \in }{ M
}
{  }{
}
{    }{
}
{   }{
}
}
{}{}{}
nur in endlich vielen der Mengen $U_j$ liegt.}
\faktzusatz {}
\faktzusatz {}

}
{

\teilbeweis {}{}{}
{Es sei die offene Überdeckung

\mathbed {W_i} {}
 {i \in I} {}
 {} {} {} {,}
gegeben. Ferner sei

\mathbed {A_n} {}
 {n \in \N} {}
 {} {} {} {,}
eine
\definitionsverweis {kompakte Ausschöpfung}{}{}
von $M$, die es nach
[[Mannigfaltigkeit/Abzählbare Basis/Kompakte Ausschöpfung/Fakt|Lemma 22.6]]
gibt. Die offenen Mengen
\mathl{A_{n+1}^{o}  \setminus A_{n-1}}{} bilden ebenfalls eine offene Überdeckung, da es zu jedem Punkt

\mavergleichskette
{\vergleichskette
{ P
}
{ \in }{ M
}
{  }{
}
{    }{
}
{   }{
}
}
{}{}{}
ein minimales

\mavergleichskette
{\vergleichskette
{ n
}
{ \in }{ \N
}
{  }{
}
{    }{
}
{   }{
}
}
{}{}{}
mit

\mavergleichskette
{\vergleichskette
{ P
}
{ \in }{ A_n
}
{  }{
}
{    }{
}
{   }{
}
}
{}{}{}
\zusatzklammer {es sei

\mavergleichskettek
{\vergleichskettek
{ A_{-1}
}
{ = }{ \emptyset
}
{  }{
}
{    }{
}
{   }{
}
}
{}{}{}} {} {}
gibt. Für dieses $n$ ist
\mathkor {} {P \not\in A_{n-1}} {und} {P \in A_n \subseteq A_{n+1} ^{o}} {.}
Indem wir die Durchschnitte
\mathl{W_i \cap \left(  A_{n+1}^{o}  \setminus A_{n-1}  \right)}{} betrachten, können wir annehmen, dass alle Mengen der Überdeckung innerhalb von einem
\mathl{A_{n+1}^{o} \setminus A_{n-1}}{} liegen.}
{}
\teilbeweis {}{}{}
{Zu jedem Punkt

\mavergleichskette
{\vergleichskette
{ P
}
{ \in }{ M
}
{  }{
}
{    }{
}
{   }{
}
}
{}{}{}
gibt es eine offene
\zusatzklammer {verträgliche} {} {}
Kartenumgebung

\mavergleichskette
{\vergleichskette
{ P
}
{ \in }{ U_P
}
{  }{
}
{    }{
}
{   }{
}
}
{}{}{,}
die in einem der $W_i$ liegt und für die es Ballumgebungen

\mavergleichskettedisp
{\vergleichskette
{ U \left(   0 , \delta_P   \right)
}
{ \subset} { B  \left(  0 , \epsilon_P  \right)
}
{ \subset} { V_P
}
{ } {
}
{ } {
}
}
{}{}{}
gibt mit

\mavergleichskette
{\vergleichskette
{ P
}
{ \in }{ \alpha_P^{-1} \left(  U \left(   0 , \delta_P   \right)  \right)
}
{  }{
}
{    }{
}
{   }{
}
}
{}{}{}
und

\mavergleichskette
{\vergleichskette
{ \delta_P
}
{ < }{ \epsilon_P
}
{  }{
}
{    }{
}
{   }{
}
}
{}{}{.}
Diese

\mathbed {\alpha_P^{-1} \left(  U \left(   0 , \delta_P   \right)  \right)} {}
 {P \in M} {}
 {} {} {} {,}
bilden dann ebenfalls eine offene Überdeckung von $M$. Nach
[[Topologischer Raum/Abzählbare Basis/Überdeckung besitzt abzählbare Teilüberdeckung/Aufgabe|Aufgabe 2.8 (Maß- und Integrationstheorie (Osnabrück 2022-2023))]]
können wir zu einer abzählbaren Teilüberdeckung davon übergehen. Wir können also annehmen, dass ein System von Karten

\mathbed {U_j} {}
 {j \in \N} {}
 {} {} {} {,}
zusammen mit Ballumgebungen

\mavergleichskettedisp
{\vergleichskette
{ U \left(   0 , \delta_j   \right)
}
{ \subset} { B  \left(  0 , \epsilon_j  \right)
}
{ \subset} { V_j
}
{ } {
}
{ } {
}
}
{}{}{}
derart gegeben ist, dass auch

\mathbed {\alpha_j^{-1} \left(  U \left(   0 , \delta_j   \right)  \right)} {}
 {j \in \N} {}
 {} {} {} {,}
eine offene Überdeckung von $M$ ist, dass jedes $U_j$ in einem
\mathl{W_{i(j)}}{} liegt und dass die oben beschriebene Beziehung zu der kompakten Ausschöpfung gilt.}
{}
\teilbeweis {}{Wir werden eine Teilmenge

\mavergleichskette
{\vergleichskette
{ J
}
{ \subseteq }{ \N
}
{  }{
}
{    }{
}
{   }{
}
}
{}{}{}
derart definieren, dass die Familie

\mathbed {U_j} {}
 {j \in J} {}
 {} {} {} {,}
auch noch die Endlichkeitseigenschaft erfüllt.\leerzeichen{}}{}
{Zu

\mavergleichskette
{\vergleichskette
{ n
}
{ \in }{ \N
}
{  }{
}
{    }{
}
{   }{
}
}
{}{}{}
betrachten wir die kompakte Menge
\mathl{A_{n+1} \setminus A_n ^{o}}{.} Diese wird von endlich vielen der

\mathbed {\alpha_j^{-1} \left(  U \left(   0 , \delta_j   \right)  \right)} {}
 {j \in \N} {}
 {} {} {} {,}
überdeckt, und zwar braucht man dazu nur Indizes $j$ mit der Eigenschaft, dass $U_j$ in
\mathl{A_{n+2}^{o}  \setminus A_{n-1}}{} liegt. Die zugehörige endliche Indexmenge sei mit $J_n$ bezeichnet, und sei

\mavergleichskette
{\vergleichskette
{ J
}
{ = }{ \bigcup_{n \in \N} J_n
}
{  }{
}
{    }{
}
{   }{
}
}
{}{}{.}
Dann wird jedes $A_n$ nur von endlich vielen der

\mathbed {U_j} {}
 {j \in J} {}
 {} {} {} {,}
getroffen.}
{}

}

__NOEDITSECTION__

\inputfaktbeweis
{Mannigfaltigkeit/Abzählbare Basis/Überdeckung/Partition/Fakt}
{Satz}
{}
{

\faktsituation {Es sei $M$ eine
\definitionsverweis {differenzierbare Mannigfaltigkeit}{}{}
mit einer
\definitionsverweis {abzählbaren Basis}{}{}
der Topologie.}
\faktfolgerung {Dann gibt es zu jeder offenen Überdeckung eine der Überdeckung untergeordnete
\definitionsverweis {stetig differenzierbare}{}{}
\definitionsverweis {Partition der Eins}{}{.}}
\faktzusatz {}
\faktzusatz {}

}
{

Nach
[[Mannigfaltigkeit/Abzählbare Basis/Lokal endlicher Atlas/Fakt|Lemma 22.9]]
können wir davon ausgehen, dass eine offene Überdeckung aus Kartengebieten

\mathbed {U_j} {}
 {j \in J} {}
 {} {} {} {,}
\zusatzklammer {$J$ abzählbar} {} {}
mit
\maabbdisp {\alpha_j} {U_j } { V_j
} {}
und mit Ballumgebungen

\mavergleichskettedisp
{\vergleichskette
{ U \left(   0 , \delta_j   \right)
}
{ \subset} { B  \left(  0 , \epsilon_j  \right)
}
{ \subset} { V_j
}
{ } {
}
{ } {
}
}
{}{}{}
\zusatzklammer {mit

\mavergleichskettek
{\vergleichskettek
{ \delta_j
}
{ < }{ \epsilon_j
}
{  }{
}
{    }{
}
{   }{
}
}
{}{}{}} {} {}
vorliegt derart, dass auch die
\mathl{\alpha_j^{-1} \left(  U \left(   0 , \delta_j   \right)  \right)}{} eine Überdeckung von $M$ bilden und dass jeder Punkt

\mavergleichskette
{\vergleichskette
{ P
}
{ \in }{ M
}
{  }{
}
{    }{
}
{   }{
}
}
{}{}{}
nur in endlich vielen der $U_j$ und insbesondere nur in endlich vielen dieser
\mathl{\alpha_j^{-1} \left(  U \left(   0 , \delta_j   \right)  \right)}{} enthalten ist. Auf $V_j$ betrachten wir die Funktion $g_j$, die durch

\mavergleichskettedisp
{\vergleichskette
{ g_j (v)
}
{ =} { \begin{cases}  e^{ - { \frac{   1  }{  \left(  \delta_j^2- \Vert { v } \Vert^2  \right)^2 } } }  \text{ für } \Vert { v } \Vert < \delta_j \, , \\ 0 \text{ sonst} \, , \end{cases}
}
{ } {
}
{ } {
}
{ } {
}
}
{}{}{}
definiert ist. Diese Funktion hat genau auf
\mathl{U \left(   0 , \delta_j   \right)}{} einen positiven Wert und ihr
\definitionsverweis {Träger}{}{}
ist
\mathl{B  \left(  0 , \delta_j  \right)}{.} Eine Betrachtung auf den beiden offenen Teilmengen
\zusatzklammer {die $V_j$ überdecken} {} {}
\mathkor {} {U \left(   0 , \epsilon_j   \right)} {und} {V_j \setminus B  \left(  0 , \delta_j  \right)} {}
zeigt, dass $g_j$ unendlich oft differenzierbar ist. Wir definieren eine Funktion
\maabbdisp {\tilde{g}_j} { M } {\R
} {}
durch

\mavergleichskettedisp
{\vergleichskette
{ \tilde{g}_j(x)
}
{ =} { \begin{cases}  g( \alpha_j(x)) \text{ für }  x \in \alpha_j^{-1}( U \left(   0 , \delta_j   \right) ) \, , \\  0 \text{ sonst} \, . \end{cases}
}
{ } {
}
{ } {
}
{ } {
}
}
{}{}{}
Diese Funktion ist
\definitionsverweis {stetig differenzierbar}{}{}
auf $M$, da der \anfuehrung{Streifen}{}
\mathl{B  \left(  0 , \epsilon_j  \right) \setminus U \left(   0 , \delta_j   \right)}{} einen glatten Übergang erlaubt. Wir setzen

\mavergleichskettedisp
{\vergleichskette
{ \tilde{g}(x)
}
{ \defeq} { \sum_{j \in J} \tilde{g}_j(x)
}
{ } {
}
{ } {
}
{ } {
}
}
{}{}{,}
wobei dies für jeden Punkt eine endliche Summe ist, da der
\definitionsverweis {Träger}{}{}
von $\tilde{g}_j$ in

\mavergleichskettedisp
{\vergleichskette
{ \alpha^{-1} \left(  B  \left(  0 , \delta_j  \right)  \right)
}
{ \subset} { U_j
}
{ } {
}
{ } {
}
{ } {
}
}
{}{}{}
liegt. Diese Funktion ist
\definitionsverweis {stetig differenzierbar}{}{}
auf $M$ und überall positiv, da die
\mathl{\tilde{g}_j(x)}{} auf den überdeckenden Mengen
\mathl{\alpha^{-1} \left(  U \left(   0 , \delta_j   \right)  \right)}{} positiv sind. Dann bilden die

\mavergleichskettedisp
{\vergleichskette
{ h_j
}
{ =} { { \frac{   \tilde{g}_j  }{ \tilde{g}  } }
}
{ } {
}
{ } {
}
{ } {
}
}
{}{}{}
die gesuchte Partition der Eins.

}

  \zwischenueberschrift{Orientierungen auf Mannigfaltigkeiten und Volumenformen}

Mit Hilfe von Partitionen der Eins können wir nun die Umkehrung von
[[Nullstellenfreie Volumenform/Impliziert orientierte (Karten) Mannigfaltigkeit/Fakt|Lemma 15.5]]
beweisen.

__NOEDITSECTION__

\inputfaktbeweis
{Nullstellenfreie Volumenform/Orientierte (Karten) Mannigfaltigkeit/Äquivalenz/Fakt}
{Satz}
{}
{

\faktsituation {Es sei $M$ eine
\definitionsverweis {differenzierbare Mannigfaltigkeit
}{}{}
mit einer
\definitionsverweis {abzählbaren Basis der Topologie}{}{.}}
\faktfolgerung {Dann existiert genau dann eine
\definitionsverweis {stetige}{}{}
nullstellenfreie
\definitionsverweis {Volumenform}{}{}
auf $M$, wenn $M$
\definitionsverweis {orientierbar}{}{}
ist.}
\faktzusatz {Diese Volumenform kann dann auch stetig differenzierbar und positiv gewählt werden.}
\faktzusatz {}

}
{

\teilbeweis {}{}{}
{Die eine Richtung wurde bereits
in [[Nullstellenfreie Volumenform/Impliziert orientierte (Karten) Mannigfaltigkeit/Fakt|Lemma 15.5]]
bewiesen.}
{}
\teilbeweis {}{}{}
{Es sei also umgekehrt $M$
\definitionsverweis {orientierbar}{}{}
und ein
\definitionsverweis {abzählbarer}{}{}
\definitionsverweis {orientierter Atlas}{}{}

\mathbed {\left(  U_i,V_i,\alpha_i  \right)} {}
 {i \in I} {}
 {} {} {} {,}
von $M$ gegeben. Dabei ist

\mavergleichskette
{\vergleichskette
{ V_i
}
{ \subseteq }{ \R^n
}
{  }{
}
{    }{
}
{   }{
}
}
{}{}{}
\definitionsverweis {offen}{}{}
und die
\definitionsverweis {Koordinaten}{}{}

\mathl{x_1 , \ldots , x_n}{} definieren eine nullstellenfreie stetige
\zusatzklammer {sogar beliebig oft differenzierbare} {} {}
\definitionsverweis {Volumenform}{}{}

\mathl{dx_1 \wedge \ldots \wedge dx_n}{} auf $V_i$. Wir setzen

\mavergleichskettedisp
{\vergleichskette
{ \omega_i
}
{ =} { \alpha_i^* dx_1 \wedge \ldots \wedge dx_n
}
{ } {
}
{ } {
}
{ } {
}
}
{}{}{}
und erhalten so eine nullstellenfreie Volumenform auf $U_i$, die wir außerhalb von $U_i$ durch $0$ fortsetzen\zusatzfussnote {Diese Fortsetzung ist natürlich nicht stetig, das spielt aber für das Folgende keine Rolle} {.} {.}

Es sei nun

\mathbed {h_j} {}
 {j \in J} {}
 {} {} {} {,}
eine der Überdeckung

\mathbed {U_i} {}
 {i \in I} {}
 {} {} {} {,}
untergeordnete, stetig differenzierbare
\definitionsverweis {Partition der Eins}{}{,}
die es
nach [[Mannigfaltigkeit/Abzählbare Basis/Überdeckung/Partition/Fakt|Satz 22.10]]
gibt. Insbesondere gibt es also für jedes $j$ ein
\mathl{i(j)}{} derart, dass der
\definitionsverweis {Träger}{}{}
von $h_j$ in
\mathl{U_{i(j)}}{} liegt. Daher sind die
\mathl{h_j \omega_{i(j)}}{}
\definitionsverweis {stetige}{}{}
$n$-Differentialformen auf $M$. Wir setzen

\mavergleichskettedisp
{\vergleichskette
{ \omega
}
{ =} { \sum_{j \in J} h_j \omega_{i(j)}
}
{ } {
}
{ } {
}
{ } {
}
}
{}{}{.}
Dies ist für jeden Punkt

\mavergleichskette
{\vergleichskette
{ P
}
{ \in }{ M
}
{  }{
}
{    }{
}
{   }{
}
}
{}{}{}
eine endliche Summe und somit eine wohldefinierte stetige $n$-Differentialform auf $M$. Für einen Punkt

\mavergleichskette
{\vergleichskette
{ P
}
{ \in }{ M
}
{  }{
}
{    }{
}
{   }{
}
}
{}{}{}
und eine die Orientierung repräsentierende
\definitionsverweis {Basis}{}{}

\mathl{v_1 , \ldots , v_n}{}  von
\mathl{T_PM}{} ist

\mavergleichskettedisp
{\vergleichskette
{ \omega(P; v_1 , \ldots , v_n)
}
{ =} { \sum_{j \in J} h_j(P) \omega_{i(j)} \left(  P; v_1 , \ldots , v_n  \right)
}
{ } {
}
{ } {
}
{ } {
}
}
{}{}{.}
Dabei gibt es ein $j$ mit

\mavergleichskette
{\vergleichskette
{ h_j(P)
}
{ > }{ 0
}
{  }{
}
{    }{
}
{   }{
}
}
{}{}{,}
und für dieses $j$ ist auch

\mavergleichskette
{\vergleichskette
{ \omega_{i(j)}(P; v_1 , \ldots , v_n)
}
{ > }{ 0
}
{  }{
}
{    }{
}
{   }{
}
}
{}{}{,}
da ja

\mavergleichskette
{\vergleichskette
{ P
}
{ \in }{ U_j
}
{  }{
}
{    }{
}
{   }{
}
}
{}{}{}
liegt, sodass diese Form überall positiv ist.}
{}

}

 [[Kategorie:Latexseite]]
